\documentclass{article}
\usepackage{graphicx} % Required for inserting images
\usepackage{amsmath}
\usepackage[margin=1in]{geometry}
\title{Model 1}
\author{Mara Dukuly and Adi Suresh}
\date{\today}

\begin{document}

\maketitle

\section{Motivation}
\begin{itemize}
    \item \textbf{Scenario}: Suppose we have a country affected by some conflict.

    \item \textbf{Assumption}: We can provide them with the money to rebuild their country after the war, and they have their pre-war levels.

    \item \textbf{Example}: The number of schools, hospitals, etc.

    \item \textbf{Observation}: War-affected countries seem to have persistent effects in terms of long-term development (i.e. Liberia, Sierra Leone, Sudan, Eritrea, Ethiopia, etc.)

    \item \textbf{Question}: What is missing?

    \item \textbf{Answer}: Human capital and skill formation.

    \item Over some period of conflict, an individual does not build their ``set'' of skills and therefore loses potential gains from skill formation they would otherwise have built during that period.

    \item We assume the individual had, or would have built, some sort of skill-capital during the war period through:
    \begin{itemize}
        \item Primary school education
        \item Parental investment
        \item Work experience
    \end{itemize}

    \item This loss then poses an inter-generational issue.
\end{itemize}

\noindent How can we aggregate individuals' conflict exposure and skill formation losses to explain current macroeconomic outcomes?

\medskip
\noindent Develop a theoretical and empirical model that captures the following:
\begin{enumerate}
    \item Capture learning-by-doing losses and skill formation as human capital in a CES production function.
    \item Measure heterogeneous effects in agents as aggregated human capital losses across multiple periods.
    % \item Estimate human-capital "catch-up" in a post-war economy. This may not be our case, but we can revise this question/ motivating point
\end{enumerate}

\section{Model}

Each time period $t$ is a 3-year period until the individual is 47. Let $a$ represent the age of the agent in the current period, where $A$ is the age in period $t_N$, the age at death. Then $a \in [0,47]$ and $t \in [0,11]$.

\medskip
\noindent The agent solves the following sequential problem:
\begin{equation}
    \max_{\{C_{a,t},\, X_{a,t},\, I_t,\, K_{t+1},\, H_{a,t+1}\}_{t=0}^{T}}
    \; \sum_{t=0}^{T} \beta^t \sum_{a=0}^{A} \frac{C_{a,t}^{1-\gamma}}{1-\gamma}
\end{equation}
subject to, for all $t = 0, \dots, T$ and $a = 0, \dots, A$:
\begin{align}
    \sum_{a=0}^{A} C_{a,t} + \sum_{a=0}^{A} X_{a,t} + I_t + K_t &\leq Y_t
        && \text{(resource constraint)} \\
    Y_t &= A K_t^{\alpha} \big(H_t L_t\big)^{1-\alpha}
        && \text{(output)} \\
    K_{t+1} &= (1-\delta) K_t + I_t
        && \text{(capital accumulation)} \\
    H_{a+1,t+1} &= \Omega_{a,t} \Big[\phi_{a,t}\, X_{a,t}^{\rho_t} + \theta_{a,t}\, H_{a,t}^{\rho_t}
        && \text{(skill formation)} \notag \\
    &\qquad + \big(1 - \phi_{a,t} - \theta_{a,t}\big)\, H_{P,t}^{\rho_t}\Big]^{1/\rho_t} \\
    \bar{H}_{t+1} &= \frac{1}{A} H_0 + \frac{A-1}{A}\, \bar{H}_t,
        \qquad H_t \equiv \bar{H}_t
        && \text{(average human capital)} \\
    C_{a,t} \geq 0, \quad & X_{a,t} \geq 0, \quad K_{t+1} \geq 0,
\end{align}
with $\rho_t \in (-\infty, 1)$, given the initial conditions $K_0$ and $\{H_{a,0}\}_{a=0}^{A}$. When agents are aged 12--15, their human capital at that point is taken as the entry level $H_0 = H_{a,t}$ in the average human capital equation.

\medskip
\noindent\textbf{Variables.}
\begin{itemize}
    \item \emph{Choice variables:} consumption $C_{a,t}$; skill investment $X_{a,t}$; physical investment $I_t$.
    \item \emph{State variables:} physical capital $K_t$; human capital of an agent of age $a$, $H_{a,t}$; average human capital $\bar{H}_t$.
    \item \emph{Other endogenous variables:} output $Y_t$; parental human capital $H_{P,t}$; entry-level human capital $H_0$.
    \item \emph{Exogenous:} labor $L_t$; TFP $A$ in the output equation.
\end{itemize}

\noindent\textbf{Parameters.}
\begin{itemize}
    \item $\beta$: discount factor.
    \item $\gamma$: coefficient of relative risk aversion.
    \item $\alpha$: capital share.
    \item $\delta$: depreciation rate of physical capital.
    \item $\Omega_{a,t}$: productivity of the skill-formation technology.
    \item $\phi_{a,t}$: weight on investment.
    \item $\theta_{a,t}$: weight on own human capital.
    \item $1-\phi_{a,t}-\theta_{a,t}$: weight on parental human capital.
    \item $\rho_t$: substitution parameter; the elasticity of substitution is $1/(1-\rho_t)$.
    \item $A$: age at death (number of ages) in the sums and in the $\bar{H}$ equation.
    \item $T$: final period.
\end{itemize}

\end{document}
